\input{../tex-inputs/latex-leseansicht-vorspann}

\section[Olga, Arthur Schnitzler und Elisabeth Gussmann an Gustav Schwarzkopf, {[}5. 6. 1908?{]}]{L04369 Olga, Arthur Schnitzler und Elisabeth Gussmann an Gustav Schwarzkopf, {[}5. 6. 1908?{]}}
\nopagebreak\mylabel{L04369v}
\rehead{ }\normalsize\beginnumbering\briefempfaengerindex{Schwarzkopf, Gustav@\textsc{Schwarzkopf, Gustav}!zzzSteinrück, Elisabeth@\emph{von Elisabeth Steinrück}!1908-06-051@{{[}5. 6. 1908?{]}}|(be}\briefempfaengerindex{Schwarzkopf, Gustav@\textsc{Schwarzkopf, Gustav}!zzzSchnitzler, Arthur@\emph{von Arthur Schnitzler}!1908-06-051@{{[}5. 6. 1908?{]}}|(be}\briefempfaengerindex{Schwarzkopf, Gustav@\textsc{Schwarzkopf, Gustav}!zzzSchnitzler, Olga@\emph{von Olga Schnitzler}!1908-06-051@{{[}5. 6. 1908?{]}}|(be}
\toendnotes[C]{\smallbreak\pagebreak[2]}
\correspDesc{Versand  durch Olga Schnitzler, Arthur Schnitzler, Elisabeth Gussmann am [5. 6. 1908?] in Hinterbrühl
\newline{}Übermittlung  am [5. 6. 1908?] in Hinterbrühl
\newline{}Erhalt  durch Gustav Schwarzkopf im Zeitraum [6. 6. 1908 – 10. 6. 1908?] in Wien}\toendnotes[C]{\smallbreak}
\Standort{DLA, A:Schnitzler, HS.1985.1.1897.}
\physDesc{Bildpostkarte, 351 Zeichen
\newline{}Handschrift Olga Schnitzler: Bleistift, lateinische Kurrent
\newline{}Handschrift Arthur Schnitzler: Bleistift, deutsche Kurrent
\newline{}Handschrift Elisabeth Steinrück: Bleistift, lateinische Kurrent
\newline{}Versand: Stempel: »\nobreak{}\oindex{Hinterbrühl@\textbf{Hinterbrühl}|pwk}Hinterbrühl, \textcolor{gray}{5}. VI. \textcolor{gray}{08}, XII\nobreak{}«.  }\toendnotes[C]{\smallbreak}\pstart{}{\pb}Herrn\pend{}\pstart{}Gustav Schwarzkopf\pend{}\pstart{}Wien I\oindex{Wien@\textbf{Wien}|pw}\pend{}\pstart{}Tiefer Graben 23.\oindex{Wien@\textbf{Wien}!I., Innere Stadt@\textbf{I., Innere Stadt}!Tiefer Graben 23@\textbf{Tiefer Graben 23}|pw}\pend{}{\bigskip}
\pstart
           {\pb}\textcolor{gray}{\textbf{HINTERBRÜHL\oindex{Hinterbrühl@\textbf{Hinterbrühl}|pw}.}}\hfill {\pb}\textcolor{gray}{\textbf{Joseph Deigner’s Hôtel Feldmarschall
                        Radetzky\oindex{Hotel Radetzky@\textbf{Hotel Radetzky}|pw}.}}\pend
           \vspace{1em}
\pstart
           \noindent{}{\pb}Lieber Herr Schwarzkopf wir sind alle heraussen, wohnen im Sanatorium Samuely\oindex{Cur- und Wasserheilanstalt Hinterbrühl@\textbf{Cur- und Wasserheilanstalt Hinterbrühl}|pw} und würden uns sehr freuen,
               Sie heraussen zu sehen, wir bleiben über \label{K_L04369-1v}\edtext{Pfingsten}{\lemma{\textnormal{\emph{Pfingsten}}}\Cendnote{\textnormal{Der Poststempel
               ist nicht mit vollständiger Sicherheit zu entziffern. Nimmt man als Orientierungspunkt, dass die drei an der Karte beteiligten Personen
               sich an Pfingsten in Hinterbrühl\oindex{Hinterbrühl@\textbf{Hinterbrühl}|pwk} aufhalten müssen, kommt nur das Jahr 1908 in Frage. In
                  diesem Jahr war der 8. 6. 1908 Pfingstsonntag.}}}\label{K_L04369-1}.\pend
           \selectlanguage{ngerman}\vspace{1em}
\pstart
           \noindent{}{\pb}{[}hs. Steinrück:{]} Lieber guter Herr Schwarzkopf, kommen Sie doch bestimmt: Ich hab Sie
               so lange nicht gesehn u. sehn mich sehr.\pend
           
\pstart
           Ihre alte {\\[\baselineskip]}\spacefill\mbox{Lisl Gussmann}\pend
           \leftskip=0em{}\selectlanguage{ngerman}\vspace{1em}\pstart {[}hs. A. Schnitzler:{]} Herzlichſt Ihr \spacefill\mbox{Arthur}.\pend{}\selectlanguage{ngerman}\endnumbering\briefempfaengerindex{Schwarzkopf, Gustav@\textsc{Schwarzkopf, Gustav}!zzzSteinrück, Elisabeth@\emph{von Elisabeth Steinrück}!1908-06-051@{{[}5. 6. 1908?{]}}|)be}\briefempfaengerindex{Schwarzkopf, Gustav@\textsc{Schwarzkopf, Gustav}!zzzSchnitzler, Arthur@\emph{von Arthur Schnitzler}!1908-06-051@{{[}5. 6. 1908?{]}}|)be}\briefempfaengerindex{Schwarzkopf, Gustav@\textsc{Schwarzkopf, Gustav}!zzzSchnitzler, Olga@\emph{von Olga Schnitzler}!1908-06-051@{{[}5. 6. 1908?{]}}|)be}\mylabel{L04369h}
\begin{anhang}
\end{anhang}\newcommand{\dateiname}{L04369}\newcommand{\titel}{Olga, Arthur Schnitzler und Elisabeth Gussmann an Gustav Schwarzkopf, [5. 6. 1908?]}\newcommand{\editorInnen}{Herausgegeben von Jahnke, SelmaMüller, Martin Anton}\input{../tex-inputs/latex-leseansicht-abspann}
